\scanpage{619}
\mspage{7}
élément de $\Phi$ est facteur direct d'un composé
transfini de flèches déduites de $\Phi_0\cup(\text{Identités de }M)$
par cochgt. de base.

Je rappelle encore les sorites $\pm$ tautologiques
des foncteurs adjoints

\uline{Proposition} 1.5. Soit

\diagram{\begin{tikzcd}[ampersand replacement=\&,column sep=large]
\text{(1.5.1)} \& M' \arrow[r,shift right,"u^{\unsure{\alpha|*}}"'] \& M \arrow[l,shift right,"u_{!}"']
\end{tikzcd}}

un couple de foncteurs adjoints, \unsure{Soient|soient} \deleted{$\Phi\unsure{\in|\subset}\mathrm{Fl}\,M$},
\deleted{partie $\mathrm{Fl}\,M$, $\Psi'$ une \unsure{part}}

\[
\Phi\subset\mathrm{Fl}(M),\qquad \Psi'\subset\mathrm{Fl}\,M'), 
\]

alors on a

\[
\big(u_!(\Phi)\longleftrightarrow\Psi'\big)\Longleftrightarrow
\big(\Phi\longleftrightarrow u^*(\Psi')\big)\tag*{(1.5.2)}
\]

\uline{Corollaire 1.6}. Pour \unsure{tte|toute} partie $\Phi$ de $\mathrm{Fl}(M)$,
\deleted{\unsure{posant} $\Psi=\Phi_*$, \unsure{on}}

\[
\big(u_!(\Phi)\big)_*=(u^*)^{-1}(\Phi_*)\tag*{(1.6.1)}
\]

En effet, appliquant 1.5. à $\Phi$ et à $\{f'\}$,
où $f'\in\mathrm{Fl}\,M$, on trouve

\[
f'\in\big(u_!(\Phi)\unsure{\big\rangle|\big)}_*\Longleftrightarrow u^*(f)\in\Phi_*,
\]

ce qui équivaut à (1.6.1). Dualement

\uline{Corollaire 1.7} : Pour tte partie $\Psi'$ de $M'$, on
a

\[
\big(u^*(\Psi')\big)^*=(u_!)^{-1}(\Psi'{}^*)\tag*{(1.7.1)}
\]

% slanted
\margin{\unsure{donner ici} 1,9. \unsure{pag|page} 9}

En effet, si $f\in\mathrm{Fl}\,M$, on a par (1.5.2)

\[
f\in\unsure{(|\{}u^*(\Psi')\unsure{\}|)}^{*}\Longleftrightarrow u_!(f)\in\Psi'{}^*,
\]
% archive sheet number: 618
